\section{H40: Agents couple per model call, not per minute}

\textbf{The question:} An agent reads a message from another agent. Is its chance to reply fixed per model call, or fixed per minute of wall time? If it is per call, then call cadence is a coupling knob, and slow agents are weakly coupled.

\textbf{The intuition:} In Glauber dynamics (a kinetic Ising model where spins update one at a time), each spin gets update attempts at some rate. The coupling per attempt is fixed, so the coupling per second is the coupling per attempt times the attempt rate. H40 maps each model call to one update attempt. This is the call clock. The rival is the wall clock: the reply chance grows with the minutes that pass, whatever the cadence. Under the call clock a slow agent is a ``heavy spin'': it has the same coupling per call but responds less per hour. The order parameter is $\eta$, the exposure-time elasticity (how fast the reply chance per call grows with the wall time that call spans). The call clock gives $\eta=0$, and the wall clock gives $\eta=1$.

\textbf{What we measured:} We fit a discrete-time reply-hazard model on each recipient's real call schedule from the context ledger. Each message enters the risk set at the recipient's read-out call (the first call whose context holds it). Replies come from the shared reply labels. We ran 33 non-reserved goal periods and four native tests. The nulls are the wall clock ($\eta=1$, refitted), a between-agent rival (slope $-1$), and a family rival (lab and style alone). For the scheduler, the risk set starts at the read-out call, and overnight read-outs are dropped. The G51 pause-timer test uses waits set before the message arrived. For the external drive, only agent-to-agent messages are items, and agent-by-unit intercepts absorb period drives. For the shared model prior, lab effects and style components enter the between-agent slope. For common input, the risk set starts after the message is in context. Synthetic replies on the real schedules came first. The estimator recovered $\eta=0$, 1 and 0.5 with bias $\le0.05$ and power 1.0 against the wrong clock.

\textbf{What we found:} In the computer-use scaffold (regimes II--III, 11 periods), coupling runs on the call clock. A call that spans more wall time carries no more reply hazard: $\eta=0.00\pm0.05$ in G51, and $-0.03$ [$-0.14$, 0.08] at timer wakes fixed in advance. The regime means are $-0.01$ [$-0.12$, 0.11] (II) and $-0.16$ [$-0.32$, 0.01] (III). In G51 the per-call coupling is flat in cadence ($s=+0.11$ [$-0.03$, 0.25]) and is a family and style trait. The per-hour coupling rises with call rate (slope 1.00 [0.81, 1.19]). The slowest fifth of agents answer within 5 minutes $8\times$ less often than the fastest fifth. In regime~I (22 periods) the call clock fails: $\eta=+0.51$ [0.41, 0.61]. Natives: G51 supported, NE41 mixed, G18 and NE14 failed. The claim that stands is the call clock in regimes II--III, with regime~I partly on the wall clock\cred{0.72}, EU 2.52. The metadata and the card agree.

\begin{figure}[h]
\includegraphics[width=\textwidth]{H40-call-clock-coupling/fig.pdf}
\caption{Agents in the computer-use scaffold couple per model call. (a) An illustration: two agents with the same reply chance per call line up when counted in calls and fan out in minutes. (b) $\eta$ per goal period: regime~I sits between the two clocks, and regimes II--III sit at the call clock (0). (c, d) In G51 the per-call coupling is flat in call rate, and the per-hour coupling rises with a slope near 1.}
\end{figure}

\textbf{What it means:} For an operator in a computer-use loop: cadence is a coupling knob that needs no prompt change. To couple an agent faster, halve its call interval. Its 5-minute reply coupling then rises by about $\times1.5$ [1.3, 1.7] in regime~III (G51 $\times1.32$), and its 30-minute coupling by about $\times1.4$. Give coordinators short pause timers. To decouple an agent, give it long batch jobs. For a scaffold builder: long tool calls, long pauses and slow models make an agent a heavy spin. For a physicist: each call behaves as one Glauber update attempt. Do not apply this rule to a scheduler-driven chat loop (regime~I), where time itself carries part of the coupling.

\textbf{Caveats:} Regime~I is not explained; endogenous generation time is untested. Agents partly choose their own cadence, and only the pause-timer dose is close to exogenous. The G51 cross-sectional slope (1.0) is larger than the model's speed-up elasticity (0.40), so composition adds to cadence. G38 and G44 have $\eta$ below zero ($-0.57$, $-0.34$), which neither clock explains. Reply labels hold one parent per message and undercount multi-target replies. The collapse test failed marginally (21/33). The clean cadence intervention (NE20) lies in reserved data. No confirmation test on reserved goal periods has run yet.
